\section*{Integrals%
\normalfont\scriptsize{ $(n\in\mathbb{N}_0,m\in\mathbb{Z},\alpha,\beta,\gamma,\delta,\mu,\nu,\sigma,r\in\mathbb{R},x\in\mathbb{R}\cap\mathrm{Dom}_f,\;a,b,z\in\mathbb{C})$}}
\underline{Basic}\\[0pt]
    \hspace*{0.2cm}
    \scalebox{0.785}{$\int (x+\alpha)^rdx \!=\!\frac{(x+\alpha)^{r+1}}{r+1}$\;\;\;\;\;$ \int x(x+\alpha)^rdx \!=\!\frac{(x+\alpha)^{r+1}(rx{+}x{-}\alpha)}{(r{+}1)(r{+}2)}$\;\;\;\;\;$\int a^xdx \!=\!\frac{a^x}{\ln a}$\;\;\;\;\;$\int u\,dv =uv{-}\int v\,du$}\\
\underline{Rational}
    \\[3pt]\hspace*{0.2cm}
        \scalebox{0.785}{$\int \frac{dx}{\alpha x+\beta} =\frac{1}{\alpha}\ln|\alpha x+\beta|$\;\;\;\; \;$\int \frac{dx}{x^2+\alpha^2} =\frac{1}{\alpha}\arctan\frac{x}{\alpha}$\;\;\;\;\;$\int \frac{dx}{x^2-\alpha^2}=\frac{1}{2\alpha}\ln |\frac{x-\alpha}{x+\alpha}|$\;\;\;\;\;$\int \frac{dx}{\alpha^2-x^2} =\frac{1}{2\alpha}\ln|\frac{\alpha{+}x}{\alpha{-}x}|$}\\[3pt]\hspace*{0.2cm}
        \scalebox{0.785}{$\int \frac{dx}{\alpha x^2+\beta x+\gamma} =\frac{2}{\surd({4\alpha\gamma{-}\beta^2})}\arctan\frac{2\alpha x{+}\beta}{\surd({4\alpha\gamma{-}\beta^2)}}$\;\;\;\;\;$\int \frac{dx}{(x{+}\alpha)(x{+}\beta)} =\frac{1}{\beta{-}\alpha}\ln|\frac{\alpha{+}x}{\beta{+}x}|$}\\[2pt]
    \underline{Roots}
    \\[3pt]\hspace*{0.2cm}
        \scalebox{0.785}{$\int\surd(x^2+\alpha^2)=\frac{1}{2}\big[x\surd(x^2+\alpha^2)+\alpha^2\operatorname{arsinh}\frac{x}{|\alpha|}\big]$\;\;\;\;\;$\int\surd(x^2-\alpha^2)=\frac{1}{2}\big[x\surd(x^2-\alpha^2)-\alpha^2\ln|\surd(x^2-\alpha^2)+x|\big]$}
    \\[3pt]\hspace*{0.2cm}
        \scalebox{0.785}{$\int \frac{dx}{\surd{(x^2+\alpha^2)}} ={\operatorname{arsinh}}\frac{x}{|\alpha|}$\;\;\;\;\;$\int \frac{dx}{\surd{(-x^2+\alpha^2)}} ={\arcsin}\frac{x}{|\alpha|}$\;\;\;\;\;$\int\frac{dx}{\surd(x^2-\alpha^2)}=\ln |\surd(x^2-\alpha^2)+x|$}
    \\[3pt]\hspace*{0.2cm}
        \scalebox{0.785}{$\int\!\frac{dx}{x\surd(x^2+\alpha^2)}\!=\!-\!\frac{1}{\alpha}\operatorname{arsinh}\frac{\alpha}{|x|} $\;\;\;\;$\int\!\frac{dx}{x\surd(-x^2+\alpha^2)}\!=\!-\frac{1}{\alpha}\!\ln\frac{|\surd (-x^2+\alpha^2) +\alpha|}{|x|}$\;\;\;\;$\int\!\frac{dx}{x\surd(x^2-\alpha^2)}\!=\! \frac{1}{\alpha}\!\arctan\frac{\surd(x^2-\alpha^2)}{\alpha}$}
    \\[3pt]\hspace*{0.2cm}
        \scalebox{0.785}{$\int \frac{x}{\surd(x^2\pm\alpha^2 )}dx=\surd(x^2\pm \alpha^2)$\;\;\;\;$\int\frac{x}{\surd(-x^2 +\alpha^2)}dx=-\surd(-x^2+\alpha^2)$}
    \\[3pt]\hspace*{0.2cm}
        \scalebox{0.785}{$\int\frac{dx}{(x^2\pm\alpha^2)^{3/2}}=\frac{\pm x}{\alpha^2 \surd(x^2\pm\alpha^2)}$\;\;\;\;$\int\frac{dx}{(-x^2+\alpha^2)^{3/2}}=\frac{x}{\alpha \surd(-x^2+\alpha^2)}$}
    \\[3pt]\hspace*{0.2cm}
        \scalebox{0.785}{$\int\frac{x}{(x^2\pm\alpha^2)^{3/2}}=\frac{-1}{\surd(x^2\pm\alpha^2)}$\;\;\;\;$\int\frac{x}{(x^2\pm\alpha^2)^{3/2}}=\frac{-1}{\surd(x^2\pm\alpha^2)}$}\\[2pt]
\underline{Trigonometric}\;\;\scalebox{0.8}{$(\mu,\nu>0, \; \chi\equiv x\alpha , \; \gamma\equiv \alpha+\beta, \; \delta \equiv \alpha-\beta)$}
     \\[3pt]\hspace*{0.2cm}
        \scalebox{0.785}{$\int \sin x\,dx = -\cos x $\;\;\;\:$\int  \cos x \, dx = \sin  x$\;\;\;\;\;\;\;\;$\int   \frac{dx}{\sin^2 x} = -\cot x $\;\;\;\;\;\;\;$\int   \frac{dx}{\cos^2 x} = \tan x $\;\;\;\;\:\;\;\:$\int \frac{dx}{\tan ^2 x} = -\cot x-x$}
    \\[3pt]\hspace*{0.2cm}
        \scalebox{0.785}{$\int \sinh x\,dx\! =\! \cosh x $\;\;\;\;\:$\int  \cosh x \, dx\!=\! \sinh x$\;\;\;\;\:$\int   \frac{dx}{\sinh^2 x}\! =\! -\coth x $\;\;\;\;\:$\int   \frac{dx}{\cosh^2 x} \!=\! \tanh x $\;\;\;\;\:$\int\frac{dx}{\tanh^2 x}\!=\!-\coth x\!+x$}
    \\[3pt]\hspace*{0.2cm}
        \scalebox{0.785}{$\int \tan x \, dx= -\ln |\cos x|$\;\;\;\;\:$\int \tan^2x\,dx =\tan x{-}x$\;\;\;\;\:$\int\tanh x\, dx=\ln \cosh x$\;\;\;\;\:$\int\tanh ^2 x\, dx=-\tanh x+x$}
    \\[5pt]\hspace*{0.2cm}
        \scalebox{0.785}{$\int \frac{dx}{\sin x} ={-}\ln|\frac{1}{\sin x}{+}\frac{1}{\tan x}|$\;\;\;\;\:\;\;\:$\int \frac{dx}{\cos x} =\ln|\frac{1}{\cos x}{+}\tan x|$\;\;\;\;\:\;\;\:$\int \frac{dx}{\tan x}=\ln|\sin x|$}
    \\[5pt]\hspace*{0.2cm}
        \scalebox{0.785}{$\int \sin^n\alpha x\,dx \!=\!{-}\frac{\sin^{n-1}\chi\,\cos\chi}{n\alpha}{+}\frac{n-1}{n}\int \sin^{n-2}\alpha x\,dx$\;\;\;\;\:$\int \cos^n\alpha x\,dx \!=\!{+}\frac{\cos^{n-1}\chi\,\sin\chi}{n\alpha}{+}\frac{n-1}{n}\int \cos^{n-2}\alpha x\,dx$}
    \\[5pt]\hspace*{0.2cm}
        \scalebox{0.785}{$\int \sin\alpha x\sin\beta x\,dx \!=\!{-}\frac{\sin \gamma x }{2\gamma}{+}\frac{\sin \delta x}{2\delta}$\;\;\;\;\:$\int \cos\alpha x\cos\beta x\,dx \!=\!{+}\frac{\sin \gamma x}{2\gamma }{+}\frac{\sin\delta x}{2\delta}$\;\;\;\;\:$\int \sin\alpha x\cos\beta x\,dx \!=\!{-}\frac{\cos\gamma x}{2\gamma}{-}\frac{\cos\delta x}{2\delta}$}
    \\[4pt]\hspace*{0.2cm}
        \scalebox{0.785}{$\int x\sin\alpha x\,dx =\frac{\sin\chi }{\alpha^2}{-}\frac{x\cos\chi}{\alpha}$\;\;\;\;\:$\int x\cos\alpha x\,dx =\frac{\cos\chi }{\alpha^2}{+}\frac{x\sin\chi}{\alpha}$\;\;\;\;\:$\int x\genfrac{}{}{0pt}{}{\sin^2}{\cos^2}\alpha x\,dx =\mp\frac{2\chi\,\sin 2\chi \, {+}\cos 2\chi \, \mp\, 2\chi^2}{8\alpha^2}$}
    \\[5pt]\hspace*{0.2cm}
        \scalebox{0.77}{$\int \!x\genfrac{}{}{0pt}{}{\sin\alpha x\sin\beta x}{\cos\alpha x\cos\beta x}dx {=}\!\mp\!\frac{x\sin \gamma x}{2\gamma}\!\mp\!\frac{\cos\gamma x}{2\gamma^2}{+}\frac{x\sin \delta x}{2\delta}{+}\frac{\cos\delta x}{2\delta^2}$\;\;\,\,$\int \!x\sin \alpha x \cos \beta x\,dx {=}{-}\frac{x\cos\gamma x}{2\gamma}{+}\frac{\sin \delta x}{2\delta^2}{-}\frac{x\cos\delta x}{2\delta}\!{+}\frac{\sin\gamma x}{2\gamma^2}$}
    \\[5pt]\hspace*{0.2cm}\scalebox{0.8}{\underline{Definite integrals} ($m!! =m(m{-}2)(m{-}4)...$, \; ${-}1!! =0!! =1!! =1$)}
    \\[2pt]\hspace*{0.2cm}
        \scalebox{0.785}{$\int_0^{\pi/2}\! \sin^\mu x\,dx {=}\int_0^{\pi/2}\! \cos^\mu x\,dx\!=\!\frac{1}{2}\!\operatorname{B}(\frac{\mu+1}{2},\frac{1}{2}) {=}\frac{(n-1)!!}{n!!}\!\cdot\!\begin{cases}\frac{\pi}{2}&\hspace{-3pt}\text{if }\mu{=}n\text{ even}\\\;1&\hspace{-3pt}\text{if }\mu {=}n\text{ odd}\end{cases}$\;\;\;\;\;\;$\int_{-1}^{+1} \genfrac{}{}{0pt}{}{\sin(m\pi x)\sin(\tilde{m}\pi x)}{\cos(m\pi x)\cos(\tilde{m}\pi x)}dx \!=\!\delta_{m,\tilde{m}}$}
    \\[5pt]\hspace*{0.2cm}
        \scalebox{0.785}
{$\left.\hspace{-5pt}
\begin{array}{c}
\int_0^\alpha\!\!\genfrac{}{}{0pt}{}{\sin(\frac{m\pi x}{\alpha})\sin(\frac{\tilde{m} \pi x}{\alpha})}{\cos(\frac{m \pi x}{\alpha})\cos(\frac{\tilde{m}\pi x}{\alpha})}dx \!=\!\frac{\alpha}{2}\delta_{m,\tilde{m}}
\vspace{3pt}\\
\text{also valid $m=1/2$}
\end{array}
\right.$\;\;\;$\int_0^\alpha \!\sin(\frac{m \pi x}{\alpha})\cos(\frac{\tilde{m}\pi x}{\alpha})dx\!=\!\begin{cases}\;\;\;\;\;0&\!\!\text{if }m{+}\tilde{m}\text{ even}\vspace{3pt}\\\frac{2m\alpha}{m^2-\tilde{m}^2}\frac{1}{\pi}&\!\!\text{if }m{+}\tilde{m}\text{ odd}\end{cases}$\;\;\;\;$\int_0^\pi\! \genfrac{}{}{0pt}{}{\sin x\,dx =2}{\cos x\,dx =0}$}
    \\[5pt]\hspace*{0.2cm}
        \scalebox{0.785}{$\int_0^\pi \sin^nx\cos^{\tilde{n}}x\,dx =0\ \forall\, \tilde{n}\text{ odd}$\;\;\;\;\;\;$\int_0^{\pi\mu} \genfrac{}{}{0pt}{}{\sin^2\alpha x}{\cos^2\alpha x}dx =\frac{1}{4\alpha}\left[2\pi\alpha\mu\mp\sin(2\pi\alpha\mu)\right]\stackrel{\text{if }\mu =n} =\frac{\pi n}{\alpha}$\;\;\;\;\;$\int_0^\pi \genfrac{}{}{0pt}{}{\sin^3x\,dx =\frac{4}{3}}{\cos^3x\,dx =0}$}
    \\[5pt]\hspace{0.2cm}
        \scalebox{0.785}{$\int_0^{2\pi} \genfrac{}{}{0pt}{}{\sin x}{\cos x}dx =0$\;\;\;\;\;\;$\int_0^{2\pi} \sin x\cos x\,dx =0$\;\;\;\;\;\;$\int_0^{2\pi} \sin^nx\cos^{\tilde{n}} x\,dx =0\text{ if }n,\tilde{n}\text{ not both even}$\;\;\;\;\;\;$\int_0^{2\pi} \genfrac{}{}{0pt}{}{\sin^3x}{\cos^3x}dx =0$}
    \\[5pt]\hspace{0.2cm}
        \scalebox{0.785}{$\int_0^{2\pi} (1{-}\cos x)^n\sin nx\,dx =0$\;\;\;\;\;\;$\int_0^{2\pi} (1{-}\cos x)^n\cos nx\,dx =(-1)^n\frac{\pi}{2^{n-1}}$}\\[5pt]
\underline{Parity}\;\;\;\;\;\;\;\scalebox{0.8}{$\text{ \textit{Even}}: f_e(-x) = f_e(x)$\;\; 
 \scalebox{0.785}{$\int_{-\alpha}^{+\alpha} f_e(x)\,dx =2\int_{0}^{\alpha}f_e(x)\,dx$\hspace{0.15cm}}\;\;\;$\text{ \textit{Odd}}: f_o(-x) = -f_o(x)$ \;\;\scalebox{0.785}{$\int_{-\alpha}^{+\alpha} f_o(x)\,dx =0$}}
    \\[5pt]\hspace{0.2cm}
        \scalebox{0.785}{$f_e  :   \cos x,\;  \cosh x,\;  x^{2n},\;  e^{-x^2} ,\;  |x|,\;  \delta_{ij},\;  \delta(x),\;  \mathbb{R}, \;  1/f_e, \;  f'_o, \; f_e\pm f_e, \;  f_e\cdot f_e, \; f_o\cdot f_o, \; \mathcal{F}\{f_e(x)\} (\xi),...$}
    \\[5pt]\hspace{0.2cm}
    \scalebox{0.785}{$f_o :  \sin x,\; \sinh x,\; x^{2n+1},\; \tan x,\; \operatorname{erf}x,\; \operatorname{sign}x,\;\ln \big(\frac{1+x}{1-x}\big),\; 1/f_o, \; f'_e,\; f_o\pm f_o, \; f_e \cdot f_o, \;\mathcal{F}\{f_o(x)\} (\xi) ,...$}\\[2pt]
\underline{Log/Exp} \; \scalebox{0.75}{$(r\neq-1)$} 
    \\[2pt]\hspace*{0.2cm}
        \scalebox{0.785}{$\int x^r\ln x\,dx =x^{r+1}\left(\frac{\ln x}{r+1}{-}\frac{1}{(r+1)^2}\right)$\;\;\;\;\;\;$\int(\ln x)^n \, dx= (-1)^n \, n!\, x \displaystyle \sum_{k=0}^n \frac{(-\ln x)^k}{k!}$\;\;\;\;\;$\int \frac{dx}{(e^{-x/\alpha}{+}1)} =\alpha\ln(e^{x/\alpha}{+}1)$}
    \\[5pt]\hspace*{0.2cm}
        \scalebox{0.785}{$\int xe^{\alpha x^2}dx =\frac{e^{\alpha x^2}}{2\alpha}$\;\;\;\;\;\;\;\;$\int x^n  e^{\alpha x}dx=e^{\alpha x}\displaystyle\sum_{k=0}^n  \frac{n!\, (-1)^k}{(n-k)!}\frac{x^{n-k}}{\alpha ^{k+1}}$\;\;\;\;\;\;\;\;$\int \frac{e^{\alpha x}}{x^n}dx =\frac{1}{n{-}1}\left({-}\frac{e^{\alpha x}}{x^{n-1}}{+}\alpha\int \frac{e^{\alpha x}}{x^{n-1}}dx\right)$}
    \\[5pt]\hspace*{0.2cm}\scalebox{0.8}{\underline{Definite integrals} ($r-1,\,\alpha>0$\;\;\;$\gamma\equiv$Euler-Mascheroni constant)}
    \\[3pt]\hspace*{0.2cm}
        \scalebox{0.785}{$\int_0^\infty x^re^{-\alpha x^2}dx =\frac{\Gamma\left(\frac{r+1}{2}\right)}{2\alpha^{\frac{r+1}{2}}} =\begin{cases}\frac{(2n{-}1)!!}{2^{n+1}\;\alpha^n}\sqrt{\frac{\pi}{\alpha}} & \text{if }r =2n \vspace{3pt}\\\frac{n!}{2\;\alpha^{n+1}} & \text{if }r =2n{+}1\end{cases}$\;\;\;\;\;}\scalebox{0.785}{\[\begin{array}{ll}&\int_0^\infty e^{-\alpha x^2}dx =\frac{1}{2}\sqrt{\frac{\pi}{\alpha}}\vspace{5pt}\\&\int_0^\infty x^2e^{-\alpha x^2}dx=\frac{1}{4}\sqrt{\frac{\pi}{\alpha^3}}\end{array}\]}
    \\[6pt]\hspace*{0.2cm}
        \scalebox{0.785}{$\int_0^\infty x^re^{-ax}dx =\frac{\Gamma(r{+}1)}{a^{r+1}}\stackrel{\text{if }r =n} =\frac{n!}{a^{n+1}}$ \ {\scriptsize$(r{>}{-}1,\Re(a){>}0)$}\;\;\;\;\;\;\;\;$\int_0^\infty \sqrt{x}e^{-x}dx =\frac{\sqrt{\pi}}{2}$\;\;\;\;\;\;\;\;$\int_0^\infty \frac{x}{e^x{-}1}dx =\frac{\pi^2}{6}$}
    \\[5pt]\hspace*{0.2cm}
        \scalebox{0.785}{$\int_0^\infty e^{-ax^b}dx =a^{-1/b}\Gamma\left(\frac{1}{b}{+}1\right)$ \;\;\;\;\;\;\;\;$\int_{-\infty}^{+\infty} e^{-\alpha x^2{+}\beta x}dx =\sqrt{\frac{\pi}{\alpha }}e^{\frac{\beta^2}{4\alpha }}$\;\;\;\;\;\;\;\;$\int_0^{2\pi} e^{i(m-\tilde{m})\phi}d\phi =2\pi\delta_{m,\tilde{m}}$}
    \\[5pt]\hspace*{0.2cm}
        \scalebox{0.8}{$\int_0^\infty e^{-\alpha x}\sin(\beta x)dx =\frac{\beta}{\alpha^2{+}b^2}$\;\;\;\;\;\;$\int_0^\infty e^{-\alpha x}\cos(\beta x)dx =\frac{\alpha }{\alpha^2{+}\beta^2}$\;\;\;\;\;\;$\int_0^\infty \frac{\ln x}{e^x}dx =\int_1^\infty \left(\frac{1}{x}{-}\frac{1}{\lfloor x\rfloor}\right)dx ={-}\gamma$}
      \\[5pt]\hspace*{0.2cm}
      \scalebox{0.8}{$\int_{0}^{\infty} x e^{-\alpha x} \sin (\beta x)
     dx = \frac{2\alpha
     \beta }{(\alpha^2 + \beta^2)^2 }$\;\;\;\;\;\;$\int_{0}^{\infty} x e^{-\alpha x} \cos (\beta x)
     dx = \frac{\alpha^2-
     \beta^2 }{(\alpha^2 + \beta^2)^2 }$\;\;\;\;\;\;$\int_{\Omega} e^{iqr \cos \theta}d\Omega = 4\pi\frac{\sin(qr)}{qr}$}
      \\[5pt]
      \hspace*{0.2cm}\scalebox{0.8}{\underline{Error function integrals} ($\varphi =\frac{1}{\sigma\sqrt{2\pi}}e^{-\frac{(x-\mu)^2}{2\sigma^2}}\!,$ \;\;$\mu{\equiv}$mean, $\sigma^2{\equiv}$variance)\;\;\;\;$\operatorname{erf}(\pm \infty) =\pm1$\;\;\;\;$i\operatorname{erfi}(z) =\operatorname{erf}(iz)$}
    \\[3pt]\hspace{0.2cm}
        \scalebox{0.785}{$\frac{2}{\sqrt{\pi}}\int_0^ze^{-t^2}dt =\operatorname{erf}(z)$\;\;\;\;\;\;$\int \varphi dx =\frac{1}{2}\operatorname{erf}\left(\frac{x{-}\mu}{\sqrt{2}\,\sigma}\right)$\;\;\;\;\;\;$\int \sqrt{x}\,e^{ax}dx =\frac{\sqrt{x}\,e^{ax}}{a}{-}\frac{\sqrt{\pi}\,\operatorname{erfi}(\sqrt{a}\sqrt{x})}{2a^{3/2}}$}\\[3pt]\hspace*{0.2cm}\scalebox{0.785}{\underline{Miscellaneous}}\\\hspace*{0.2cm}
        \scalebox{0.785}{$\tfrac{d|x|}{dx} = \mathrm{sgn}(x)$\;\;\; $\tfrac{d^2 |x|}{d^2x}=\delta (x)$\;\;\;$\langle \sin x\rangle = \langle \cos x\rangle=\sqrt{2}/2\;\;\; \langle \sin^2 x\rangle = \langle\cos ^2 x \rangle = 1/2$ }


\newcolumn  

\section{Integrales típicas en QED}

$$\int_{-1}^{1} \left( 1 + \cos^2\theta \right) d(\cos\theta) = \frac{8}{3}  \quad 
\quad \int_{-1}^{1} \frac{d(\cos\theta)}{(1 - \beta\cos\theta)^2} = \frac{2}{1-\beta^2}$$



$$\int \frac{A + B\cos\theta + C\cos^2\theta}{(1 + \cos\theta)^2} d(\cos\theta)\! =\!  [x\! =\! \cos \theta, u \!=\! 1+x ] \!=\! \int \left( C + \frac{B - 2C}{u} + \frac{A - B + C}{u^2} \right) du $$\\
$$= Cu + (B - 2C)\log|u| - \frac{A - B + C}{u} + K \Rightarrow C(1+x) + (B - 2C)\log|1+x| - \frac{A - B + C}{1+x} + K$$\\
$$[K' = K + C]  \Rightarrow  I = C\cos\theta + (B - 2C)\log|1+\cos\theta| - \frac{A - B + C}{1+\cos\theta} + K'$$


Para salvar la divergencia en $\cos \theta =-1$, se integra de $-1 + \epsilon $ a $1- \epsilon$. La retrodivergencia indica que la partícula rebota exactamente en la dirección opuesta a su trayectoria incidente en el sistema centro de masas. Señala una resonancia cinemática: la probabilidad del proceso se dispara porque la partícula virtual intercambiada se acerca a su capa de masas ($p^2 \to 0$), volviéndose ``casi real''.





$$\int \frac{A + B\cos\theta + C\cos^2\theta}{(1 - \cos\theta)^2} d(\cos\theta)\! =\! [x\! =\! \cos \theta, u \!=\! x-1 ] \!=\! \int \left( C + \frac{B + 2C}{u} + \frac{A + B + C}{u^2} \right) du $$

$$= Cu + (B + 2C)\log|u| - \frac{A + B + C}{u} + K \Rightarrow C(x-1) + (B + 2C)\log|x-1| - \frac{A + B + C}{x-1} + K$$

$$[K' = K - C] \Rightarrow I = C\cos\theta + (B + 2C)\log|1-\cos\theta| + \frac{A + B + C}{1-\cos\theta} + K'$$


Para salvar la divergencia en $\cos \theta = 1$, se integra de $-1 + \epsilon$ a $1 - \epsilon$. La dispersión hacia adelante (forward scattering) indica que la partícula apenas se desvía de su trayectoria incidente en el sistema centro de masas. Señala una resonancia cinemática en el canal $t$: la probabilidad del proceso se dispara porque la partícula virtual intercambiada se acerca a su capa de masas ($t \to 0$), volviéndose ``casi real''. Se puede entender, desde el punto de vista clásico, como que las partículas interactúan a una distancia infinita sin desviarse de su trayectoria original. Ocurre aún cuando $m\neq 0$. 

\color{cyan} ¡Estás integrando con $d(\cos\theta) \Rightarrow$ variable a sustituir según $\lim$s integración es $\cos\theta$, no $\theta$! \color{black}

$$\int(\frac{1+x}{1-x} + \frac{1-x}{1+x}) \, dx = \int \frac{1+x^2}{1-x^2}\, dx = -x+2 \,\mathrm{ arctanh} (x) \quad [\text{Integrando par y } \rm arctanh(0)=0]$$  



$$\int \frac{(1+x)^2}{(1-x)^2}dx= x - \frac{4}{x - 1} + 4 \ln\left(\left|x - 1\right|\right)\qquad \int \frac{\left( 1 - x\right)^{2}}{\left( x + 1\right)^{2}} dx = -4 \ln\left(\left| x + 1 \right|\right) - \frac{4}{x + 1} + x$$\\ $$ \int \frac{1}{(1-x)^2}dx=\frac{1}{1 - x}\;\;\; \int \frac{1}{\left(x + 1\right)^{2}}dx = -\frac{1}{x + 1}$$


\hrule



\section*{Desintegraciones Exactas con Masas Finitas}

\textbf{Quark Top a Bosón W Real ($t \to qW^+$)}


La anchura de desintegración exacta a nivel árbol (sin límite $m_q \to 0$):
$$\Gamma(t \to qW^+) = \frac{G_F m_t^3}{8\pi\sqrt{2}} |V_{tq}|^2 \lambda^{1/2}(1, x_W, x_q) \left[ (1-x_q)^2 + x_W(1+x_q) - 2x_W^2 \right]$$

Donde:

  - $x_i = m_i^2/m_t^2$ son las fracciones de masa al cuadrado.
  
- $\lambda(a,b,c) = a^2 + b^2 + c^2 - 2ab - 2ac - 2bc$ es la función cinemática de Källén.


\textbf{Bosón de Higgs a Bosones Vectoriales ($H \to VV$)}

Para bosones vectoriales masivos ($V = W^\pm, Z^0$) con masa $m_V < m_H/2$ y $x_V = m_V^2/m_H^2$:
$$\Gamma(H \to W^+W^-) = \frac{G_F m_H^3}{8\pi\sqrt{2}} \sqrt{1 - 4x_W} (1 - 4x_W + 12x_W^2)$$
$$\Gamma(H \to ZZ) = \frac{G_F m_H^3}{16\pi\sqrt{2}} \sqrt{1 - 4x_Z} (1 - 4x_Z + 12x_Z^2)$$

\textit{Nota:} El factor diferencial $1/2$ en el canal $ZZ$ proviene de los bosones finales idénticos.


\section*{Fenomenología de Interacciones Fuertes y Simetrías}

\textbf{Escalado del Espacio de Fases (Branching Ratios)}

Para calcular la relación entre anchuras parciales $\Gamma_i$ asumiendo simetría de isospín exacta, la dependencia cinemática está dominada por el momento en el centro de masas $p_c$ y el momento angular orbital relativo $L$ del estado final:
$$\Gamma \propto |\mathcal{M}|^2 \, p_c^{2L+1}$$

- \textbf{Onda S ($L=0$):} Típico en $0^+ \to 0^- + 0^-$. Escala linealmente: $\Gamma \propto p_c$.
  
    - \textbf{Onda P ($L=1$):} Típico en $1^- \to 0^- + 0^-$. Escala al cubo: $\Gamma \propto p_c^3$.


Cálculo de $p_c$ para $M \to m_1 + m_2$: $p_c = \frac{\lambda^{1/2}(M^2, m_1^2, m_2^2)}{2M}$.



\textbf{Estadística de Bose-Einstein (Estados Finales Idénticos)}

Un estado final de dos bosones idénticos (e.g., $\pi^0\pi^0$) requiere una función de onda total \textbf{estrictamente simétrica} ante el intercambio de las partículas.

  
  - Simetría Espacial: Dictada por $(-1)^L$.
 
  
  - Simetría de Isospín: Acoplo de dos isospines idénticos $I_1 \otimes I_1 = I_{tot}$. 


- \textbf{Regla Crítica:} Para un estado $\pi\pi$ ($1 \otimes 1 \to 0 \oplus 1 \oplus 2$), los subestados con $I_{tot} = 0, 2$ son simétricos, mientras que $I_{tot} = 1$ es antisimétrico. Si una resonancia con $I=1$ y paridad tal que $L$ es par decae a dos piones neutros, \textbf{está rigurosamente prohibida}.


\section*{Física de Colisionadores Hadrónicos}

\textbf{Topologías Tevatron ($p\bar{p}$) vs LHC ($pp$)}

Diferencia fundamental en la composición de los Partones (PDFs) incidentes:

    
    - \textbf{Tevatron ($p\bar{p}$):} Interacciones como $u\bar{d} \to \dots$ o $q\bar{q} \to \dots$ proceden dominantemente mediante \textbf{quarks de valencia}. Ambos partones tienen fracciones de momento de Bjorken altas ($x \sim 1/3$).
    
    
    - \textbf{LHC ($pp$):} Cualquier proceso de dispersión aniquilación que requiera un antiquark ($\bar{q}$) inicial, exige que éste provenga del \textbf{mar de quarks (sea quarks)}. La PDF asociada $f_{\bar{q}/p}(x)$ cae exponencialmente con $x$, suprimiendo fuertemente la sección eficaz a altas transferencias de momento en comparación con procesos inducidos por gluones ($gg \to \dots$).


% \section*{Extra}

% \textbf{Supresión de Helicidad en el Pion ($\pi^- \to \ell^- \bar{\nu}_\ell$)}

% Aunque cinemáticamente el decaimiento al electrón ($m_e \approx 0.5$ MeV) posee mucho más espacio de fases que al muón ($m_\mu \approx 105.7$ MeV), el pion decae al muón el $99.99\%$ de las veces.

    
%     - \textbf{Causa:} El pion tiene espín 0. Por conservación de momento angular, el leptón y el antineutrino (helicidad $+\frac{1}{2}$ estricta) deben emitirse con la misma quiralidad.
    
%      - El acoplo V-A de la interacción débil fuerza al fermión masivo a adoptar la quiralidad "equivocada" (diestra). Esta penalización es proporcional a $m_\ell^2$. Como $m_\mu \gg m_e$, el canal muónico es el dominante.
    
%     - $\frac{\Gamma(\pi \to e\nu)}{\Gamma(\pi \to \mu\nu)} = \left(\frac{m_e}{m_\mu}\right)^2 \left( \frac{m_\pi^2 - m_e^2}{m_\pi^2 - m_\mu^2} \right)^2 \approx 1.2 \times 10^{-4}$.


% \textbf{Paridad G en Sistemas Multipion}


% La Paridad G ($G = C e^{i\pi I_2}$) es una simetría conservada estrictamente por la fuerza fuerte.
% Para un sistema de $n$ piones, la paridad G es \textbf{multiplicativa}:
% $$G(n\pi) = (-1)^n$$
% \textbf{Aplicación directa:} 

% - Mesón $\rho$ ($I^G = 1^+$): Decae a $2\pi$ (permitido, $(-1)^2 = +1$), pero no a $3\pi$ (prohibido, $(-1)^3 = -1$).
 
% - Mesón $\omega$ ($I^G = 0^-$): Decae a $3\pi$ (permitido), pero no a $2\pi$ (prohibido).


% \textbf{Narrow Width Approximation en el Pico de Resonancia}

% Para un scattering partónico $ab \to X \to cd$ a energías de centro de masas exactamente en el polo de la resonancia ($\sqrt{s} = m_X$), la sección eficaz se simplifica analíticamente sin integración usando Branching Ratios:
% $$\sigma_{peak}(ab \to cd) \approx \frac{16\pi}{m_X^2} \, \frac{(2J_X + 1)}{(2S_a + 1)(2S_b + 1)} \, \frac{C_X}{C_a C_b} \, \text{BR}(X \to ab) \, \text{BR}(X \to cd)$$
% Donde $S_i$ son los espines y $C_i$ son los grados de libertad de color de las partículas.

\newcolumn


\section*{Coordinate Systems}
\underline{Spherical} $(\theta \in [0, \pi], \phi \in [0, 2\pi))$
\begin{align*}
&\begin{cases} 
x = r \sin\theta\cos\phi \\ 
y = r \sin\theta\sin\phi \\ 
z = r \cos\theta 
\end{cases}
&
\begin{cases} 
\mathbf{\hat{x}} = \sin\theta\cos\phi\,\mathbf{\hat{r}} + \cos\theta\cos\phi\,\bm{\hat{\theta}} - \sin\phi\,\bm{\hat{\phi}} \\
\mathbf{\hat{y}} = \sin\theta\sin\phi\,\mathbf{\hat{r}} + \cos\theta\sin\phi\,\bm{\hat{\theta}} + \cos\phi\,\bm{\hat{\phi}} \\
\mathbf{\hat{z}} = \cos\theta\,\mathbf{\hat{r}} - \sin\theta\,\bm{\hat{\theta}} 
\end{cases}
\\[0pt]
&\begin{cases} 
r = \sqrt{x^2+y^2+z^2} \\ 
\theta = \arctan(\sqrt{x^2+y^2}/z) \\ 
\phi = \arctan(y/x) 
\end{cases}
&
\begin{cases} 
\mathbf{\hat{r}} = \sin\theta\cos\phi\,\mathbf{\hat{x}} + \sin\theta\sin\phi\,\mathbf{\hat{y}} + \cos\theta\,\mathbf{\hat{z}} \\ 
\bm{\hat{\theta}} = \cos\theta\cos\phi\,\mathbf{\hat{x}} + \cos\theta\sin\phi\,\mathbf{\hat{y}} - \sin\theta\,\mathbf{\hat{z}} \\ 
\bm{\hat{\phi}} = -\sin\phi\,\mathbf{\hat{x}} + \cos\phi\,\mathbf{\hat{y}} 
\end{cases}
\end{align*}

\underline{Cylindrical}  $(\rho \in [0, \infty), \phi \in [0, 2\pi))$
\begin{align*}
&\begin{cases} 
x = \rho\cos\phi \\ 
y = \rho\sin\phi \\ 
z = z 
\end{cases} & 
\begin{cases} 
\mathbf{\hat{x}} = \cos\phi\,\bm{\hat{\rho}} - \sin\phi\,\bm{\hat{\phi}} \\ 
\mathbf{\hat{y}} = \sin\phi\,\bm{\hat{\rho}} + \cos\phi\,\bm{\hat{\phi}} \ \ \
\\ 
\mathbf{\hat{z}} = \mathbf{\hat{z}} 
\end{cases}
\\[0pt]
&\begin{cases} 
\rho = \sqrt{x^2+y^2} \\ 
\phi = \arctan(y/x) \\ 
z = z 
\end{cases}
&
\begin{cases} 
\bm{\hat{\rho}} = \cos\phi\,\mathbf{\hat{x}} + \sin\phi\,\mathbf{\hat{y}} \\ 
\bm{\hat{\phi}} = -\sin\phi\,\mathbf{\hat{x}} + \cos\phi\,\mathbf{\hat{y}} \\ 
\mathbf{\hat{z}} = \mathbf{\hat{z}} 
\end{cases}
\end{align*}



\hrule

{

\centering

\centering
\renewcommand{\arraystretch}{1.2}
\begin{tabular}{ll}
\toprule
\textbf{Mesón (Masa en MeV)} & $\mathbf{I^G(J^{PC})}$ \\ 
\midrule
\multicolumn{2}{c}{\textbf{Ligeros sin sabor neto} ($u, d, s$)} \\
\midrule
$\pi(140)$ & $1^-(0^{-+})$ \\
$\eta(548), \eta'(958)$ & $0^+(0^{-+})$ \\
$\rho(770)$ & $1^+(1^{--})$ \\
$\omega(782), \phi(1020)$ & $0^-(1^{--})$ \\
$a_0(980)$ & $1^-(0^{++})$ \\
$f_0(980)$ & $0^+(0^{++})$ \\
$a_1(1260)$ & $1^-(1^{++})$ \\
$f_1(1285)$ & $0^+(1^{++})$ \\
$b_1(1235)$ & $1^+(1^{+-})$ \\
$h_1(1170)$ & $0^-(1^{+-})$ \\
$a_2(1320)$ & $1^-(2^{++})$ \\
$f_2(1270), f_2'(1525)$ & $0^+(2^{++})$ \\
$\pi_2(1670)$ & $1^-(2^{-+})$ \\
$\eta_2(1645)$ & $0^+(2^{-+})$ \\
$\rho_3(1690)$ & $1^+(3^{--})$ \\
$\omega_3(1670), \phi_3(1850)$ & $0^-(3^{--})$ \\
$a_4(2040)$ & $1^-(4^{++})$ \\
$f_4(2050)$ & $0^+(4^{++})$ \\
\midrule
\multicolumn{2}{c}{\textbf{Charmonium} ($c\bar{c}$)} \\
\midrule
$\eta_c(2984)$ & $0^+(0^{-+})$ \\
$J/\psi(3097), \psi(3686)$ & $0^-(1^{--})$ \\
$\chi_{c0}(3415)$ & $0^+(0^{++})$ \\
$\chi_{c1}(3510)$ & $0^+(1^{++})$ \\
$h_c(3525)$ & $0^-(1^{+-})$ \\
$\chi_{c2}(3556)$ & $0^+(2^{++})$ \\
\midrule
\multicolumn{2}{c}{\textbf{Bottomonium} ($b\bar{b}$)} \\
\midrule
$\eta_b(9399)$ & $0^+(0^{-+})$ \\
$\Upsilon(9460)$ & $0^-(1^{--})$ \\
$\chi_{b0}(9859)$ & $0^+(0^{++})$ \\
$\chi_{b1}(9893)$ & $0^+(1^{++})$ \\
$h_b(9899)$ & $0^-(1^{+-})$ \\
$\chi_{b2}(9912)$ & $0^+(2^{++})$ \\
\midrule
\multicolumn{2}{c}{\textbf{Sabor abierto} ($K, D, B$) -- Solo $\mathbf{I(J^P)}$} \\
\midrule
$K(494), D(1865), B(5279)$ & $1/2(0^-)$ \\
$K^*(892), D^*(2010), B^*(5325)$ & $1/2(1^-)$ \\
$D_s(1968), B_s(5367)$ & $0(0^-)$ \\
$D_s^*(2112), B_s^*(5415)$ & $0(1^-)$ \\
\bottomrule
\end{tabular}

}


% =====================================================================
% DERIVACIÓN GENERAL DE PDFs TIPO DELTA (MODELO ESTÁTICO DE QUARKS)
% =====================================================================
\section*{Derivación General de PDFs de Valencia}

Sea un hadrón $H$ definido exclusivamente por su contenido de valencia, compuesto por un número total de $N$ quarks constituyentes. Definimos $n_q$ como el número neto de quarks de un sabor específico $q$ presentes en el hadrón ($\sum_{q} n_q = N$).

La función de distribución: 
$$f_q(x) = n_q \delta\left(x - \frac{1}{N}\right)$$


\section*{LIPS para calcular Anchuras $\Gamma$}

Para el caso general de dos partículas finales con masas distintas $m_1 \neq m_2$ en el sistema de referencia del centro de masas (donde la energía total es $E_{\text{CM}} = M$), el módulo del trimomento final $|\mathbf{p}_f|$ viene determinado por la función de Källén:$|\mathbf{p}_f| = \frac{\sqrt{\lambda(M^2, m_1^2, m_2^2)}}{2M}$

Así, integrando sobre $d\Omega$: 
$$\int d\text{LIPS}_2 = \frac{1}{8\pi} \sqrt{1 - 2\frac{m_1^2 + m_2^2}{M^2} + \frac{(m_1^2 - m_2^2)^2}{M^4}}$$

Ojo, asegura siempre que $M > m_1 + m_2$

% =====================================================================

\clearpage


    
ÚLTIMAS PINCELADAS: 

Al calcular $|\mathcal{M}|^2$ (o $F_C$), aplica $\dagger$ y luego \underline{debes} cambiar los índices correspondientes a índices mudos (todo lo que no sea índices de las patas externas). Los índices asociados a patas externas deben mantenerse fijos.

$$\text{Tip Casimir}: \quad  \sum_{s_1, s_2} \left[ \bar{v}_1 \Gamma_1 u_2 \right] \left[ \bar{u}_2 \Gamma_2 v_1 \right] = \text{Tr} \left[ \cancel{p}_1 \Gamma_1 \cancel{p}_2  \Gamma_2 \right]$$


Para un producto: 

$$\sum_{\text{spin}}
\underbrace{[\bar{u}_1 \,\Gamma_1\, u_2]}_{\text{escalar}}
\underbrace{[\bar{u}_3 \,\Gamma_2\, u_4]}_{\text{escalar}}
\underbrace{[\bar{u}_2 \,\bar{\Gamma}_1\, u_1]}_{\text{escalar}}
\underbrace{[\bar{u}_4 \,\bar{\Gamma}_2\, u_3]}_{\text{escalar}} =$$


$$ = \sum_{\text{spin}}
[\bar{u}_1 \,\Gamma_1\, u_2][\bar{u}_2 \,\bar{\Gamma}_1\, u_1]
\times
[\bar{u}_3 \,\Gamma_2\, u_4][\bar{u}_4 \,\bar{\Gamma}_2\, u_3]=$$
$$ =\text{Tr}\!\left[\Gamma_1\,(\not{p}_2 + m_2)\,\bar{\Gamma}_1\,(\not{p}_1 + m_1)\right]
\times
\text{Tr}\!\left[\Gamma_2\,(\not{p}_4 + m_4)\,\bar{\Gamma}_2\,(\not{p}_3 + m_3)\right]$$






\vspace{2pt}
En caso U.R., en el CM siempre se cumple $|\vec{p}_i| = |\vec{p}_f|$


\vspace{2pt}

Al calcular la sección eficaz, si las dos partículas finales son indistinguibles (la misma), se multiplica por $1/2$ para no doblecontar.

\vspace{2pt}


Calcula cada peso de color (el $F_C$ de un canal particular o interferencia de ellos): 

$$F_{ab} = \frac{1}{N_{in}} \sum C_a C_b^{*}$$
$\bullet$ No se suman los distintos ``pesos'' de color entre sí! 

$\bullet$ No necesariamente igual interferencia a-b que b-a

$\bullet$ Si la estructura del factor de color es igual en varios canales $\rightarrow$ Mismo $F_C$ para ellos (incluso sus interferencias)





\vspace{2pt}

$$(a^\mu b^\nu)(c_\mu d_\nu) = (a\cdot c) (b\cdot d ) \qquad g_{\mu \nu} a^\mu b^\nu = a\cdot b$$
$$(a g^{\mu\nu})(bg_{\mu\nu}) =ab g^{\mu \nu}g_{\mu \nu} =a b \cdot D\quad \text{D}\equiv\text{dim espacio}$$ 


\begin{center}
    
\boxed{
$$S^{\mu\nu}(k_1, k_2) = k_1^\mu k_2^\nu + k_1^\nu k_2^\mu - g^{\mu\nu}(k_1 \cdot k_2)$$
}

\boxed{
$$S^{\mu\nu}(p_1, p_2) S_{\mu\nu}(p_3, p_4) = 2 \left[ (p_1 \cdot p_3)(p_2 \cdot p_4) + (p_1 \cdot p_4)(p_2 \cdot p_3) \right]$$
}


 \boxed{$$\varepsilon^{\mu\nu\rho\sigma} p_{i\rho} p_{j\sigma} \varepsilon_{\mu\nu\alpha\beta} p_k^\alpha p_l^\beta = -2 \left[ (p_i \cdot p_k)(p_j \cdot p_l) - (p_i \cdot p_l)(p_j \cdot p_k) \right]$$}
\end{center}

\vfill

\vspace{2pt}



$$\renewcommand{\arraystretch}{1.75}
\begin{array}{|l|}
\hline
\multicolumn{1}{|c|}{\textbf{EW: Álgebra Quiral y Proyectores}} \\
\hline
P_{L,R} = \tfrac{1\mp\gamma^5}{2},\quad P_{L,R}^2 = P_{L,R},\quad P_L P_R = 0 \\
P_L + P_R = 1,\quad \text{Tr}[P_{L,R}] = 2 \\
\gamma^\mu P_{L,R} = P_{R,L}\gamma^\mu,\quad P_{L,R}\cancel{a} = \cancel{a}P_{R,L} \\
\hline
\multicolumn{1}{|c|}{\textbf{EW: Identidades de Contracción}} \\
\hline
\gamma^\mu \gamma^5 \gamma_\mu = -4\gamma^5 \\
\gamma^\mu(1\mp\gamma^5)\gamma_\mu = 4(1\pm\gamma^5) \\
\gamma^\mu(1\mp\gamma^5)\cancel{a}\gamma_\mu = -2\cancel{a}(1\mp\gamma^5) \\
\gamma^\mu(1\mp\gamma^5)\cancel{a}\cancel{b}\gamma_\mu = 4(a\cdot b)(1\pm\gamma^5) \\
\gamma^\mu(1\mp\gamma^5)\cancel{a}\cancel{b}\cancel{c}\gamma_\mu = -2\cancel{c}\cancel{b}\cancel{a}(1\mp\gamma^5) \\
\hline
\multicolumn{1}{|c|}{\textbf{EW: Trazas de Dirac}} \\
\hline
\text{Tr}\big[\gamma^\mu\gamma^\nu\gamma^\rho\gamma^\sigma\big] = 4(g^{\mu\nu}g^{\rho\sigma} - g^{\mu\rho}g^{\nu\sigma} + g^{\mu\sigma}g^{\nu\rho}) \\
\text{Tr}\big[\gamma^\mu\gamma^\nu\gamma^\rho\gamma^\sigma\gamma^5\big] = 4i\varepsilon^{\mu\nu\rho\sigma} \\
\hline
\text{Tr}\big[\gamma^\mu(1\mp\gamma^5)\cancel{p}_1\gamma^\nu(1\mp\gamma^5)\cancel{p}_2\big] = \\
\quad 8\big[p_1^\mu p_2^\nu + p_2^\mu p_1^\nu - g^{\mu\nu}(p_1{\cdot}p_2) \mp i\varepsilon^{\mu\alpha\nu\beta}p_{1\alpha}p_{2\beta}\big] \\
\hline
\text{Tr}\big[\gamma^\mu(1\mp\gamma^5)\cancel{a}\gamma^\nu(1\pm\gamma^5)\cancel{b}\big] = \\
\quad 8\big[a^\mu b^\nu + b^\mu a^\nu - g^{\mu\nu}(a{\cdot}b) \pm i\varepsilon^{\mu\alpha\nu\beta}a_{\alpha}b_{\beta}\big] \\
\hline
\text{Tr}\big[\gamma^\mu(1-\gamma^5)\cancel{p}_1\gamma^\nu(1+\gamma^5)\cancel{p}_2\big] = 0 \quad (\text{ortog.}) \\
\hline
\text{Tr}\big[\gamma^\mu(1-\gamma^5)\cancel{p}_1 (1+\gamma^5)\gamma^\nu\cancel{p}_2\big] = 8\big[p_1^\mu p_2^\nu + p_2^\mu p_1^\nu - g^{\mu\nu}(p_1{\cdot}p_2) + i\varepsilon^{\mu\alpha\nu\beta}(p_1)_{\alpha}(p_2)_{\beta}\big] \\
\hline
\text{Tr}\big[\gamma^\mu(1-\gamma^5)\cancel{p}_1 (1-\gamma^5)\gamma^\nu\cancel{p}_2\big] = 0 \\
\hline
\text{Tr}\big[\cancel{p}_1\gamma_\mu(g_L P_L + g_R P_R)\cancel{p}_2(g_L P_R + g_R P_L)\gamma_\nu\big] =\\
\quad \text{Tr}\big[\cancel{p}_1\gamma_\mu \cancel{p}_2(g_L P_L + g_R P_R)^2\gamma_\nu\big] =\\
\quad 2(g_L^2+g_R^2)\big[p_{1\mu}p_{2\nu}+p_{2\mu}p_{1\nu}-g_{\mu\nu}(p_1{\cdot}p_2)\big] +\\
\quad 2i(g_R^2-g_L^2)\varepsilon^{\rho\mu\sigma\nu}p_{1\rho}p_{2\sigma} \\
\hline
\text{Tr}\big[\gamma^\mu(g_V - g_A\gamma^5)\cancel{p}_1\gamma^\nu(g_V-g_A\gamma^5)\cancel{p}_2\big] = \\
\quad 4(g_V^2+g_A^2)\big[p_1^\mu p_2^\nu+p_2^\mu p_1^\nu-g^{\mu\nu}(p_1{\cdot}p_2)\big] \\
\quad - 8ig_V g_A\varepsilon^{\mu\alpha\nu\beta}p_{1\alpha}p_{2\beta} \\
\hline
\text{Tr}\big[\gamma^\mu(g_V-g_A\gamma^5)\gamma^\nu(g_V-g_A\gamma^5)\big] = 4(g_V^2-g_A^2)g^{\mu\nu} \\
\hline
\text{Tr}\big[(\cancel{p}_1+m_1)\gamma^\mu(g_V-g_A\gamma^5)(\cancel{p}_2-m_2)\gamma^\nu(g_V-g_A\gamma^5)\big] = \\
\quad 4(g_V^2+g_A^2)\big[p_1^\mu p_2^\nu+p_2^\mu p_1^\nu-g^{\mu\nu}(p_1{\cdot}p_2)\big] \\
\quad - 8ig_Vg_A\varepsilon^{\mu\alpha\nu\beta}p_{1\alpha}p_{2\beta} - 4m_1m_2(g_V^2-g_A^2)g^{\mu\nu} \\
\hline
\end{array}$$

\newcolumn



\begin{center}
    PROPIEDAD CÍCLICA TRAZA: $\operatorname{Tr}\{\cancel{a}\gamma^\mu\cancel{b}\gamma^\nu\} = \operatorname{Tr}\{\gamma^\nu \cancel{a}\gamma^\mu\cancel{b}\}$
\end{center}

$$\renewcommand{\arraystretch}{1.5}
\begin{array}{|l|l|}
\hline
\textbf{Expresión / Traza} & \textbf{Resultado} \\
\hline
\multicolumn{2}{|c|}{\textbf{QED: Identidades de Contracción}} \\
\hline
\gamma^\mu \gamma_\mu & 4 \\
\gamma^\mu \cancel{a} \gamma_\mu & -2\cancel{a} \\
\gamma^\mu \cancel{a} \cancel{b} \gamma_\mu & 4(a \cdot b) \\
\gamma^\mu \cancel{a} \cancel{b} \cancel{c} \gamma_\mu & -2\cancel{c}\cancel{b}\cancel{a} \\
\hline
\multicolumn{2}{|c|}{\textbf{QED: Trazas de Dirac}} \\
\hline
\text{Tr}[\mathbf{1}] & 4 \\
\text{Tr}[\gamma^{\mu_1} \dots \gamma^{\mu_{2n+1}}] & 0 \\
\text{Tr}[\cancel{a}\cancel{b}] & 4(a \cdot b) \\
\text{Tr}[\cancel{a}\cancel{b}\cancel{c}\cancel{d}] & 4[(a \cdot b)(c \cdot d) - (a \cdot c)(b \cdot d) + (a \cdot d)(b \cdot c)] \\
\text{Tr}[\gamma^5] & 0 \\
\text{Tr}[\gamma^5 \cancel{a}\cancel{b}] & 0 \\
\text{Tr}[\gamma^5 \cancel{a}\cancel{b}\cancel{c}\cancel{d}] & -4i\epsilon_{\mu\nu\rho\sigma} a^\mu b^\nu c^\rho d^\sigma \\
\text{Tr}[\gamma_\mu \cancel{a} \gamma^\mu \cancel{b}] & -8(a \cdot b) \\
\text{Tr}[\gamma_\mu \cancel{a} \cancel{b} \gamma^\mu \cancel{c} \cancel{d}] & 16(a \cdot b)(c \cdot d) \\
\text{Tr}[\gamma_\mu (\cancel{a} - \cancel{b}) \cancel{a} (\cancel{a} - \cancel{b}) \gamma^\mu \cancel{c}] & -16(a \cdot b)(b \cdot c) \quad (a^2=b^2=0) \\
\text{Tr}[\cancel{c} \gamma^\mu (\cancel{a} - \cancel{b}) \gamma^\nu \cancel{a} \gamma_{\color{magenta}\nu \color{black}}(\cancel{a} - \cancel{\color{magenta}b\color{black}}) \gamma_{\color{magenta}\mu\color{black}}] & 32(a \cdot b)(b \cdot c) \quad (a^2=b^2=0) \\
\text{Tr}[\cancel{c} \gamma^\mu (\cancel{a} - \cancel{b}) \gamma^\nu \cancel{a} \gamma_{\color{magenta}\nu\color{black}} (\cancel{a} - \cancel{\color{magenta}d\color{black}}) \gamma_{\color{magenta}\mu\color{black}}] & 16 [ (a \cdot b)(c \cdot d) - (a \cdot c)(b \cdot d) + (a \cdot d)(b \cdot c)] \\\text{Tr}[\cancel{c} \gamma^\mu (\cancel{a} - \cancel{b}) \gamma^\nu \cancel{a} \gamma_{\color{magenta}\mu\color{black}} (\cancel{a} - \cancel{\color{magenta}b\color{black}}) \gamma_{\color{magenta}\nu\color{black}}] &  64(a \cdot b)(a \cdot c) \quad (a^2=b^2=0)\\
\text{Tr}[\cancel{c} \gamma^\mu (\cancel{a} - \cancel{b}) \gamma^\nu \cancel{a} \gamma_{\color{magenta}\mu\color{black}} (\cancel{a} - \cancel{\color{magenta}d\color{black}}) \gamma_{\color{magenta}\nu\color{black}}] &   32(a \cdot c)(a \cdot b + a \cdot d - b \cdot d)\quad (a^2=b^2=0)\\

\hline
\multicolumn{2}{|c|}{\textbf{QCD: Álgebra de Color } SU(3) \textbf{ } (t^a = \lambda^a/2)} \\
\hline
\text{Tr}[t^a] & 0 \\
\text{Tr}[t^a t^b] & T_R \delta^{ab} \quad (T_R = 1/2) \\
\text{Tr}[t^a t^b t^c] & \frac{1}{4} (d^{abc} + i f^{abc}) \\
\text{Tr}[t^a t^b t^a t^c] & -\frac{1}{12} \delta^{bc} \\
t^a t^a & C_F \mathbf{1}_{3\times3} \quad (C_F = 4/3) \\
f^{acd} f^{bcd} & C_A \delta^{ab} \quad (C_A = 3) \\
\text{Tr}[t^a t^b] \text{Tr}[t^a t^c] & \frac{1}{4} \delta^{bc} \\
\text{Tr}[t^a t^b t^c] \text{Tr}[t^a t^b t^d] & \frac{2}{3} \delta^{cd} \\
\hline
\end{array}$$




\begin{center}
    $$(\bar{\psi}_1 \Gamma \psi_2)^\dagger = \bar{\psi}_2 \overline{\Gamma} \psi_1\qquad \bar{\bar{\Gamma}} =\Gamma$$
\end{center}


$$\renewcommand{\arraystretch}{1.8}
\begin{array}{|l|l|l|}
\hline
\textbf{Operador } \Gamma & \textbf{Estructura / Corriente} & \textbf{Adjunto } \bar{\Gamma} \\
\hline
I & \text{Escalar } (S) & I \\
\hline
\gamma^\mu & \text{Vectorial } (V) & \gamma^\mu \\
\hline
\sigma^{\mu\nu} = \frac{i}{2}[\gamma^\mu, \gamma^\nu] & \text{Tensorial } (T) & \sigma^{\mu\nu} \\
\hline
\gamma^\mu \gamma^5 & \text{Axial } (A) & \gamma^\mu \gamma^5 \\
\hline
\gamma^5 & \text{Pseudoescalar } (P) & -\gamma^5 \\
\hline
\cancel{p} = p_\mu \gamma^\mu & \text{Momento (con } p_\mu \in \mathbb{R} \text{)} & \cancel{p} \\
\hline
P_R = \frac{1 + \gamma^5}{2} & \text{Proyector Quiralidad Right} & P_L = \frac{1 - \gamma^5}{2} \\
\hline
P_L = \frac{1 - \gamma^5}{2} & \text{Proyector Quiralidad Left} & P_R = \frac{1 + \gamma^5}{2} \\
\hline
\gamma^\mu (g_L^e P_L + g^e_R P_R)&\text{Comb. lineal de proyectores}& (g_L^e P_R + g^e_R P_L)\gamma^\mu\\
\hline
\gamma^\mu (1 \mp \gamma^5) & \text{Vértice V$\mp$A (Corriente Débil)} & \gamma^\mu (1 \mp \gamma^5) \\
\hline
\cancel{a}\cancel{b} & \text{Producto par de momentos} & \cancel{b}\cancel{a} \\
\hline
\cancel{a}\cancel{b}\cancel{c} & \text{Producto impar de momentos} & \cancel{c}\cancel{b}\cancel{a} \\
\hline
\gamma^\mu \cancel{a} \gamma^\nu & \text{Vértice generalizado} & \gamma^\nu \cancel{a} \gamma^\mu \\
\hline
\gamma^\mu (\cancel{p}_1 + m) \gamma^\nu & \text{Propagador flanqueado} & \gamma^\nu (\cancel{p}_1 + m) \gamma^\mu \\
\hline
c \Gamma & \text{Operador con constante } c \in \mathbb{C} & c^* \bar{\Gamma} \\
\hline
\Gamma_1 \Gamma_2 \dots \Gamma_n & \text{Cadena general} & \bar{\Gamma}_n \dots \bar{\Gamma}_2 \bar{\Gamma}_1 \\
\hline
\end{array}$$


\begin{center}
    
ÍNDICES DE COLOR EN QCD:
\setlength{\tabcolsep}{3pt}
\renewcommand{\arraystretch}{1.}
\begin{tabular}{|p{0.9cm}|p{1.3cm}|p{2cm}|p{3.5cm}|}
\hline
\textbf{Elem.} & \textbf{Cuándo} & \textbf{Dónde} & \textbf{Índice y Criterio} \\ \hline
\textbf{Ext.} & Reales (in/out) & Extremo libre & $q$: Fund. ($i \in \{1,2,3\}$) \newline $g$: Adj. ($a \in \{1,\dots,8\}$) \\ \hline
\textbf{Int.} & Virtuales (prop.) & Ambos extremos & $q$: Fundamentales ($i, j$) \newline $g$: Adjuntos ($a, b$) \\ \hline
\textbf{Nodo} & Vértices QCD & Hereda de líneas (No olvidar índices griegos de vértice) & $q$: $T^a_{ij}$ (Saliente $\rightarrow$ Entrante) \newline $g$: $f^{abc}$ (Antihorario) \\ \hline
\end{tabular}

\end{center}

 \newcolumn


{
\centering


\begin{tabular}{cccc}
\toprule
\textbf{Partícula} & \textbf{Símbolo} & \textbf{Tipo de campo} & \textbf{Propagador Explícito} \\ \midrule
Fotón & $\gamma$ & Vectorial ($M=0$) & $\displaystyle \frac{-i g_{\mu\nu}}{k^2 + i\epsilon}$ \\ \addlinespace[0.3cm]
Bosón $Z$ & $Z^0$ & Vectorial masivo & $\displaystyle \frac{-i \left(g_{\mu\nu} - \frac{k_\mu k_\nu}{M_Z^2}\right)}{k^2 - M_Z^2 + iM_Z\Gamma_Z}$ \\ \addlinespace[0.3cm]
Bosón $W$ & $W^\pm$ & Vectorial masivo & $\displaystyle \frac{-i \left(g_{\mu\nu} - \frac{k_\mu k_\nu}{M_W^2}\right)}{k^2 - M_W^2 + iM_W\Gamma_W}$ \\ \addlinespace[0.3cm]
Bosón de Higgs & $h^0$ & Escalar masivo & $\displaystyle \frac{i}{k^2 - M_h^2 + iM_h\Gamma_h}$ \\ \bottomrule
\end{tabular}

}



\begin{itemize}
    \item \textbf{\color{red}Desaparición de la anchura en el propagador SI EN CANAL $t$ Y $u$ ($\Gamma_V \to 0$):} 
    \[ \frac{1}{t - M_V^2} \quad \text{ó} \quad \frac{1}{u - M_V^2} \]

    \item \textbf{Si término $W$ o $Z$ como pata externa, introduces $\epsilon$. Relación importante:}
    $$\sum_{\lambda=1}^3 \varepsilon_Z^\mu(p, \lambda) \varepsilon_Z^{*\nu}(p, \lambda) = -g^{\mu\nu} + \frac{p_Z^\mu p_Z^\nu}{M_Z^2}$$


    \item \textbf{Límite de bajas energías (Teoría de Fermi):} A energías muy inferiores a la masa del bosón ($|t|, |u| \ll M_V^2$), el momento de la partícula virtual es despreciable frente a su masa. El propagador pierde su dependencia dinámica con el momento y colapsa en una interacción de contacto constante:
    \[ \frac{-i g_{\mu\nu}}{t - M_V^2} \xrightarrow{|t| \ll M_V^2} \frac{i g_{\mu\nu}}{M_V^2} \]
    Este es el fundamento del acoplamiento efectivo de la teoría de Fermi para corrientes cargadas, donde se define $\frac{G_F}{\sqrt{2}} = \frac{g_w^2}{8M_W^2}$.

    \item \textbf{Límite ultrarelativista ($m_f \to 0$):} Al contraer los propagadores vectoriales masivos con corrientes fermiónicas conservadas, el término quiral $\propto k_\mu k_\nu / M_V^2$ es proporcional a la masa del fermión externo por la ecuación de Dirac ($k_\mu J^\mu \propto m_f$). En el límite de altas energías ($E \gg m_f$), este término se anula estrictamente, reduciendo el numerador a $-i g_{\mu\nu}$.
    \item \textbf{Prescripción de Feynman vs. Anchura de Breit-Wigner:} El término $+i\epsilon$ en el fotón desplaza analíticamente los polos fuera del contorno de integración real. Para los bosones masivos, la sustitución $M^2 \to M^2 - iM\Gamma$ (obtenida mediante la resumación Dyson de la auto-energía a todos los órdenes) es obligatoria únicamente en procesos de tipo canal $s$ donde la energía en el centro de masas $\sqrt{s}$ se aproxima a la masa de la resonancia. En canales $t$ y $u$, dado que $k^2 < 0$, el polo físico nunca se alcanza y se puede aproximar $\Gamma \to 0$ manteniendo la prescripción $+i\epsilon$.
\end{itemize}

Si tengo neutrinos...

\begin{itemize}
    \item \textbf{Anulación de $k_\mu k_\nu$ (Exacta):} Dado que $m_\nu \to 0$, el término $\frac{k_\mu k_\nu}{M_V^2} \to 0$ de forma exacta. El numerador del propagador masivo se reduce a $-ig_{\mu\nu}$.

    \item \textbf{Acoplamientos quirales ($Z^0\nu\bar{\nu}$):} Al ser el neutrino neutro ($Q=0$) y con isospín $T^3 = \frac{1}{2}$:
    \[ g_R = -Q \sin^2\theta_W = 0, \quad g_L = T^3 - Q \sin^2\theta_W = \frac{1}{2} \]
    El vértice de corriente neutra se reduce a una estructura puramente levógira (left-handed):
    \[ \frac{-ig_w}{2\cos\theta_W} \gamma^\mu P_L \]
    
    \item \textbf{Base V-A pura:} Por la ausencia de carga ($Q=0$), los acoplamientos vectorial y axial se igualan:
    \[ c_V = T^3 - 2Q\sin^2\theta_W = \frac{1}{2}, \quad c_A = T^3 = \frac{1}{2} \]
    La corriente neutra toma la forma operacional $\gamma^\mu(1-\gamma^5)$, idéntica a la interacción cargada del $W^\pm$, maximizando la simplificación del álgebra de trazas de Dirac.


    $$\text{\textbf{Contracción Trazas}: }$$
    
    $$(\text{Sim} + \text{Antisim}) \times (\text{Sim} + \text{Antisim}) \longrightarrow (\text{Sim} \times \text{Sim}) + (\text{Antisim} \times \text{Antisim})$$


$$\varepsilon^{\mu\nu\alpha\beta} \varepsilon_{\mu\nu\rho\sigma} = -2(g^\alpha_\rho g^\beta_\sigma - g^\alpha_\sigma g^\beta_\rho)$$
% \item \textbf{Acoplamientos quirales ($Z^0\nu\bar{\nu}$):} Al ser el neutrino neutro ($Q=0$) y con isospín $T^3 = \frac{1}{2}$:
    % \[ g_R = -Q \sin^2\theta_W = 0, \quad g_L = T^3 - Q \sin^2\theta_W = \frac{1}{2} \]
    % El vértice de corriente neutra se reduce a una estructura puramente levógira (left-handed):
    % \[ \frac{-ig_w}{2\cos\theta_W} \gamma^\mu P_L \]
    
    % \item \textbf{Base V-A pura:} Por la ausencia de carga ($Q=0$), los acoplamientos vectorial y axial se igualan:
    % \[ c_V = T^3 - 2Q\sin^2\theta_W = \frac{1}{2}, \quad c_A = T^3 = \frac{1}{2} \]
    % La corriente neutra toma la forma operacional $\gamma^\mu(1-\gamma^5)$, idéntica a la interacción cargada del $W^\pm$, maximizando la simplificación del álgebra de trazas de Dirac.
\end{itemize}




